\originalsection{18.315：作业3（原卷 Fall 2005）}
收录原题3、4, 以及原题5的路径多面体衔接题, 共五个任务单元.
原题2的交错排列数在题3解答中作为必要定义与辅助计数说明,
不将其一般计数组合学题面另计为已收录图论题.

\begin{exercise}\label{mit:18315-hw3-q3}
\textbf{18.315, 原卷 Fall 2005, 作业3, 题3.}
记 $T_n(x,y)$ 为完全图 $K_n$ 的 Tutte 多项式,
$u_n$ 为满足 $\sigma_1<\sigma_2>\sigma_3<\cdots$ 的 $[n]$ 排列数.
证明 $u_n=|T_{n+1}(1,-1)|$.
\end{exercise}
\sourceof{官方 HW3 第1页题3; 编者独立解答, 下述证明还得出该取值非负}
\begin{solution}
先求交错排列的指数生成函数, 约定 $u_0=u_1=1$.
对长度 $n+1\ge2$ 的交错排列, 最大元只能位于偶数位置;
若左边有 $j$ 个元素, 则 $j$ 为奇数. 选择左边的标签后,
左右两段各有 $u_j,u_{n-j}$ 种交错次序; 反向的交错模式用标签取补便具有同一计数.
由最小元所在的奇数位置同样得到偶数 $j$ 的分解. 将这两种分别等于 $u_{n+1}$ 的计数相加:
\[
2u_{n+1}=\sum_{j=0}^{n}\binom nj u_j u_{n-j}\qquad(n\ge1).
\]
所以 $U(z)=\sum_{n\ge0}u_nz^n/n!$ 满足
$U'=(U^2+1)/2$, $U(0)=1$.
$\sec z+\tan z$ 满足同一方程和初值, 逐项比较系数有唯一解,
故 $U(z)=\sec z+\tan z$.

再给每个图赋权 $w^{|E|}$, 标签顶点数为 $r$ 时, 全部图的总权重为
$(1+w)^{\binom r2}$. 记 $C_w(z)$ 为连通图总权重的指数生成函数.
每个图唯一分成无序的连通分量, 按分量的标签划分展开指数级数, 得
\[
\exp C_w(z)=\sum_{r\ge0}(1+w)^{\binom r2}\frac{z^r}{r!}.
\]
这是一条形式幂级数恒等式, 不涉及实分析收敛问题.
由 Tutte 多项式的子集定义, 对 $r\ge1$ 有
\[
\sum_{\substack{A\subseteq E(K_r)\\([r],A)\text{ 连通}}}w^{|A|}
=w^{r-1}T_r(1,1+w).
\]
令 $w=-2$. 因为 $(-1)^{\binom r2}$ 的系数序列依次为 $1,1,-1,-1,\ldots$,
得到 $C_{-2}(z)=\log(\cos z+\sin z)$, 并且
\begin{align*}
C_{-2}'(z)
&=\sum_{n\ge0}(-2)^nT_{n+1}(1,-1)\frac{z^n}{n!}\\
&=\frac{\cos z-\sin z}{\cos z+\sin z}
=\sec(2z)-\tan(2z)=U(-2z).
\end{align*}
比较系数即 $T_{n+1}(1,-1)=u_n$. 因 $u_n$ 是非负整数,
原题带绝对值的等式随之成立.
\end{solution}

\begin{exercise}\label{mit:18315-hw3-q4}
\textbf{18.315, 原卷 Fall 2005, 作业3, 题4.}
将 $K_n$ 的生成树以顶点1为根.
若 $i<j$ 且 $j$ 在从 $i$ 到根1的路上, 称 $(i,j)$ 为一个反序对,
记树 $t$ 的反序数为 $\operatorname{inv}(t)$.
将 $f_n(q)=\sum_t q^{\operatorname{inv}(t)}$ 用 $T_n(1,y)$ 表示.
\end{exercise}
\sourceof{官方 HW3 第1页题4; 编者独立解答, 给出可逆的连通子图分割}
\begin{solution}
结论是 $f_n(q)=T_n(1,q)$. 为证明它, 对任意连通子图 $A$ 从根1作深度优先搜索,
每次优先选择标签最大的尚未访问邻点, 得到一棵根树 $t$.
无向 DFS 的非树边只能连接祖先与后代:
若连接两个不同分支, 首先被访问的端点就会把另一个分支提前搜索进去, 产生矛盾.

设非树边连接祖先 $a$ 与后代 $i$, 从 $a$ 向 $i$ 的树路的第一个子顶点为 $j$.
在搜索首次从 $a$ 进入 $j$ 时, $i$ 尚未访问.
为使最大标签规则仍选择 $j$, 必须 $i<j$.
反过来, 给 $t$ 加任意一些形如 $\{\operatorname{parent}(j),i\}$ 的边,
其中 $j$ 是 $i$ 的真祖先且 $i<j$, 搜索仍得到 $t$:
在一个顶点的尚未搜索子树中, 每条新增边的后代端点标签小于它所属分支的子根,
因此最大候选总是最大的未搜索子根;
递归搜索完该分支后, 相关新增端点已经访问, 不再改变后续选择.
按子树递归即可验证全过程.

这样, 得到树 $t$ 的全部连通子图恰为 $t\cup S$,
$S$ 任取上述允许的非树边子集. 允许边与树反序对一一对应,
且不同对不产生同一边: 对祖先 $a$ 和后代 $i$, 树路上的子根 $j$ 唯一.
因此
\[
\sum_{A\text{ 连通}}w^{|A|-n+1}
=\sum_t(1+w)^{\operatorname{inv}(t)}.
\]
左边正是 $T_n(1,1+w)$, 将 $q=1+w$ 代入便得到答案.
例如 $n=3$ 的三棵树反序数为 $0,0,1$, 所以 $f_3(q)=2+q$,
与 $T_3(1,q)=1+1+q$ 相符.
\end{solution}

\begin{exercise}\label{mit:18315-hw3-q5}
\textbf{18.315, 原卷 Fall 2005, 作业3, 题5.}
设
\[
P_n=\{x\in\R^n:x_i\ge0,\quad x_i+x_{i+1}\le1\ (1\le i<n)\}.
\]
(a) 求其整数点数; (b) 求其体积;
(c) 对正整数 $k$, 给出 $kP_n$ 的整数点数的组合解释, 并与(b)联系.
\end{exercise}
\sourceof{官方 HW3 第1页题5; 以路径稳定集多面体作为衔接选题; 原题 $n=1$ 无界, 下述有限结论明确取 $n\ge2$}
\begin{solution}
当 $n=1$ 时只有 $x_1\ge0$, 整数点和体积均无限;
原题没有写出的 $n\ge2$ 条件不可略去. 以下每个坐标至少出现在一个相邻约束中,
故 $0\le x_i\le1$.

(a) 整数点各坐标为0或1, 相邻坐标不能都为1,
即路径 $P_n$ 的独立集特征向量. 设其数为 $a_n$;
末坐标为0时有 $a_{n-1}$ 种, 为1时倒数第二坐标必须为0, 有 $a_{n-2}$ 种.
以组合计数的空路径、单顶点路径为初值 $a_0=1,a_1=2$,
得到 $a_n=F_{n+2}$, 其中 $F_0=0,F_1=1$.
这里 $a_1$ 是递推使用的路径计数, 不等同于原无界集合 $P_1$ 的整数点数.

(b) 作保持体积的仿射变换: 奇数 $i$ 取 $y_i=x_i$,
偶数 $i$ 取 $y_i=1-x_i$.
约束成为单位立方体中的
$y_1\le y_2\ge y_3\le y_4\ge\cdots$.
忽略坐标相等的零体积边界, 立方体按坐标严格大小次序分成 $n!$ 个等体积单纯形,
每个体积为 $1/n!$. 合法的坐标排名恰为 $u_n$ 个交错排列,
所以 $\operatorname{vol}(P_n)=u_n/n!$.

(c) 对 $kP_n$ 作 $y_i=x_i$（奇数 $i$）、$y_i=k-x_i$（偶数 $i$）,
其整数点恰是取值于 $\{0,1,\ldots,k\}$、满足同一交错弱不等式的整数序列.
也可将 $1<2>3<4>\cdots$ 看作一个篱笆偏序:
这些序列就是向该整数链的保序映射.
对应地, 对每个 $j=0,\ldots,k-1$ 取理想
$I_j=\{i:y_i\le j\}$, 得到 $k$ 个递增、允许相等的序理想;
反过来由一个元素首次进入理想的阶段恢复 $y_i$, 因而这是双射.
与(b)的联系是, 每种严格排名恰有 $\binom{k+1}{n}=k^n/n!+O(k^{n-1})$ 个整数序列,
重复坐标所在的有限个超平面只贡献 $O(k^{n-1})$;
因此整数点数为 $u_n k^n/n!+O(k^{n-1})$, 首项系数正是该体积.
\end{solution}
